\documentclass[a4paper]{article}
% Español (México) con XeLaTeX: fontspec para fuentes Unicode, polyglossia
% para la división silábica y los nombres del documento en la variante mexicana.
\usepackage{fontspec}
\usepackage{polyglossia}
\setdefaultlanguage[variant=mexican]{spanish}
% Los nombres chinos van como «pinyin (original)»: xeCJK da a los caracteres
% Han su propia fuente; sin ella XeLaTeX los omite con un aviso.
% FandolSong no cubre algunos caracteres de nombres (U+73FA); WenQuanYi Zen Hei sí.
\usepackage[AutoFallBack=true]{xeCJK}
\setCJKmainfont{FandolSong-Regular.otf}
\setCJKfallbackfamilyfont{\CJKrmdefault}{WenQuanYi Zen Hei}
% polyglossia no traduce los nombres de listings.
\AtBeginDocument{\providecommand{\lstlistingname}{}\renewcommand{\lstlistingname}{Listado}%
  \providecommand{\lstlistlistingname}{}\renewcommand{\lstlistlistingname}{Índice de listados}}
\usepackage{amsmath, amssymb}
\usepackage{graphicx}
\usepackage[margin=2.5cm]{geometry}
\usepackage[most]{tcolorbox}
\usepackage{etoolbox}
\usepackage{listings}
\usepackage{hyperref}
\usepackage{booktabs}
\usepackage{subcaption}
\usepackage{float}

\newtcolorbox{knowledgebox}[1]{enhanced, colback=blue!5!white, colframe=blue!75!black, colbacktitle=blue!75!black, coltitle=white, fonttitle=\bfseries, title={#1}, attach boxed title to top left={yshift=-2mm, xshift=2mm}, boxrule=1pt, sharp corners}
\newtcolorbox{importantbox}[1]{enhanced, colback=yellow!10!white, colframe=yellow!80!black, colbacktitle=yellow!80!black, coltitle=black, fonttitle=\bfseries, title={#1}, sharp corners}
\newtcolorbox{warningbox}[1]{enhanced, colback=red!5!white, colframe=red!75!black, colbacktitle=red!75!black, coltitle=white, fonttitle=\bfseries, title={#1}, sharp corners}

\lstset{language=Python, basicstyle=\ttfamily\small, keywordstyle=\color{blue}, stringstyle=\color{red!60!black}, commentstyle=\color{green!60!black}, breaklines=true, frame=single, numbers=left, numberstyle=\tiny\color{gray}, captionpos=b, extendedchars=true}

\newcommand{\notetitle}{[LLM Agents SP25] Inference-Time Techniques for LLM Reasoning — Xinyun Chen}
\newcommand{\noteauthors}{Elaboradas a partir de la clase de Xinyun Chen}
\newcommand{\notedate}{2025-01-27}
\newcommand{\videochannel}{Berkeley RDI}
\newcommand{\videopublishdate}{2025-01-27}
\newcommand{\videoduration}{1:21:32}
\newcommand{\videourl}{https://www.youtube.com/watch?v=g0Dwtf3BH-0}
\newcommand{\videocoverpath}{cover.jpg}
\newcommand{\repourl}{https://github.com/hqhq1025/ai-course-notes}

\usepackage{fancyhdr}
\pagestyle{fancy}
\fancyhf{}
\fancyfoot[C]{\thepage}
\fancyfoot[R]{\footnotesize\color{gray}\href{\repourl}{AI Course Notes}}
\renewcommand{\headrulewidth}{0pt}
\renewcommand{\footrulewidth}{0pt}

\begin{document}

\begin{titlepage}
\centering
{\Large Notas del curso\par}
\vspace{1.2cm}
{\huge\bfseries \notetitle\par}
\vspace{0.8cm}
{\large \noteauthors\par}
\vspace{0.3cm}
{\large \notedate\par}
\vspace{1.2cm}

\ifdefempty{\videocoverpath}{
{\small Al generar las notas, indique la ruta local de la portada del video.\par}
}{
\includegraphics[width=0.82\textwidth,height=0.45\textheight,keepaspectratio]{\videocoverpath}\par
}

\vfill
\begin{tcolorbox}[width=0.9\textwidth, colback=black!2!white, colframe=black!60, sharp corners]
\textbf{Autor o canal del video}: \videochannel\par
\textbf{Fecha de publicación}: \videopublishdate\par
\textbf{Duración del video}: \videoduration\par
\ifdefempty{\videourl}{}{\textbf{Enlace del video}: \href{\videourl}{\nolinkurl{\videourl}}\par}
\textbf{Página del proyecto}: \href{\repourl}{\nolinkurl{\repourl}}\par
\end{tcolorbox}
\end{titlepage}

\tableofcontents
\newpage

\section{Introducción}
Esta clase comienza con Xinyun Chen mencionando que «\textit{2025 will be the year of Agents}», y explicando que este semestre el curso Advanced LLM Agents SP25 de Berkeley RDI, en colaboración con Google DeepMind y Cal, cuenta tanto con casos consolidados en el track de aplicaciones como con un nuevo track de investigación centrado en inference-time reasoning. El docente enfatiza que más de 15,000 inscritos al MOOC, más de 3,000 participantes en el hackathon y más de 200,000 dólares en recursos de la industria buscan mostrar que este curso no es una demostración puramente teórica, sino que de verdad pretende implementar las técnicas de tiempo de inferencia en escenarios de ingeniería complejos.

\begin{knowledgebox}{Posicionamiento del curso}
De manera modular, se divide la técnica de tiempo de inferencia en tres partes: «prompt + búsqueda + iteración». El docente compartirá en esta clase, con detalle, la estrategia, las plantillas de instrucciones y los puntos clave de ingeniería de cada una de ellas, de modo que, bajo el tema de «invertir cómputo en el tiempo de inferencia», atendamos a la vez las tres dimensiones de tokens, ancho y profundidad.
\end{knowledgebox}

\subsection{Resumen de la sección}
La introducción establece el contexto de arranque en frío de los Agentes y enfatiza el inference-time compute como una palanca que mejora el razonamiento sin modificar los pesos, organizando el contenido en tres partes: prompt, búsqueda e iteración; las secciones siguientes se desarrollan conforme a esta lógica.
\section{Prompting en tiempo de inferencia y presupuesto de tokens}

\subsection{Ampliar el presupuesto de tokens para exponer la cadena de razonamiento}
Chen usa el rompecabezas de Arc AGI como ejemplo para mostrar la curva del presupuesto de tokens en tiempo de inferencia: los modelos tradicionales se detienen por debajo del 25\% con un budget de aproximadamente 20 dólares, mientras que OpenAI O3, con un presupuesto de entre 20 y 1,000 dólares y acompañado de prompts de ejemplo y un trace de ejecución, puede llevar la exactitud hasta el 87.5\%. Esta diapositiva (figura \ref{fig:token-budget}) muestra con claridad la tendencia no lineal en la parte media de la curva de precisión de razonamiento a medida que se alarga el presupuesto de tokens.

\begin{figure}[H]
\centering
\includegraphics[width=0.92\textwidth]{frame-token-budget.jpg}
\caption{La curva en tiempo de inferencia de Arc AGI y la disposición de costo-token de OpenAI O3.\protect\footnotemark}
\label{fig:token-budget}
\end{figure}
\footnotetext{Intervalo de tiempo en el video: 00:08:56--00:10:20, usa la curva de O3 para mostrar que la expansión de tokens puede alcanzar una exactitud a nivel humano.}

\begin{importantbox}{Tres lecciones de la curva en tiempo de inferencia}
1) Cuando la complejidad de la tarea aumenta, el presupuesto de tokens debe pasar de decenas a miles para ver una mejora significativa; 2) aumentar el presupuesto debe acompañarse de prompts de ejemplo y de una puntuación por pasos (stepwise scoring), de lo contrario un prompt más largo por sí solo no eleva la exactitud; 3) el modo de presupuesto alto es adecuado para tareas de nivel de competencia o de verificación, no para conversaciones nativas de baja latencia.
\end{importantbox}

\subsection{Restricciones del diseño de prompts}
El prompt necesita dos capacidades: primero, alentar al modelo a generar una Chain-of-Thought larga, para que surja de manera natural una secuencia de razonamiento coherente a lo largo de varias rondas; segundo, incorporar en el prompt una señal de dificultad de la tarea, de modo que el modelo ajuste automáticamente la profundidad de generación. En el minuto 44:31, Chen señala explícitamente: «el prompt solo puede producir un razonamiento más largo si el modelo al que se le da el prompt sabe que enfrenta un problema difícil».

\begin{knowledgebox}{Dos criterios del diseño de prompts}
1. CoT largo: usar plantillas por pasos, pares de pregunta y respuesta o demostraciones con varios ejemplos, para que el modelo dé seguimiento continuo a su estado actual; 2. marca de dificultad: mostrar en el prompt la dificultad del problema (problem difficulty), la longitud de contexto (context length) y la profundidad esperada (expected depth), para que el modelo ajuste de manera adaptativa la profundidad en tiempo de inferencia.
\end{knowledgebox}

\subsection{Observabilidad de costos y control de presupuesto}
Ampliar el presupuesto de tokens debe tomar como referencia el costo real: la diferencia entre 1 dólar, 20 dólares y 1,000 dólares por pregunta no es solo cuestión de dinero, también implica latencia y consumo de cómputo. El riesgo real expuesto es que «un presupuesto alto vuelve la tarea inexplicable y también puede consumir los recursos de los microservicios».

\begin{warningbox}{Riesgo de que el presupuesto se descontrole}
En cuanto el presupuesto en tiempo de inferencia se eleva a miles de dólares, es necesario limitarlo a subtareas de alto valor; de lo contrario, la frecuencia de llamadas queda fuera de sincronía, el rendimiento del sistema cae e incluso puede provocar una pérdida de equilibrio. Se recomienda definir de antemano una estrategia de caché de muestreo (sampling) y reservar la búsqueda (search) más costosa para la ruta crítica verificada (verified-critical path).
\end{warningbox}

\subsection{Resumen de la sección}
El prompting en tiempo de inferencia necesita encontrar un equilibrio entre «alentar un CoT largo» y «percibir la dificultad de la tarea», al mismo tiempo que se incorporan a la supervisión el presupuesto de tokens, el costo y las métricas de latencia. La gráfica de la curva de Arc AGI nos recuerda que ampliar la cantidad de tokens es eficaz, pero el costo se satura de manera gradual, por lo que el prompt, la verificación y el presupuesto deben operar de forma coordinada.

\section{Programación del cómputo en tiempo de inferencia}

\subsection{Programación triangular Token/Width/Iteration}
Chen, en el minuto 15:32, usa una diapositiva que dibuja «prompt → búsqueda → iteración» como un triángulo de token / width / iteration, y enfatiza que no basta con agregar tokens, sino que también hay que organizar la búsqueda con múltiples candidatos y varias rondas de autocorrección (figura \ref{fig:compute-pipeline}). En la práctica de ingeniería, lo habitual es usar primero el prompt para que el modelo genere candidatos más largos, después controlar el branching de cada step, y por último reservar una o dos rondas de feedback loop para cada candidate.

\begin{figure}[H]
\centering
\includegraphics[width=0.92\textwidth]{frame-pipeline.jpg}
\caption{El pipeline de tres etapas del cómputo en tiempo de inferencia: prompt → search → iteration.\protect\footnotemark}
\label{fig:compute-pipeline}
\end{figure}
\footnotetext{Intervalo de tiempo en el video: 00:15:10--00:16:10, muestra la curva de programación de token-width-iteration.}

\begin{knowledgebox}{Cuantificación de ingeniería de la programación triangular}
1. Token: medir el token budget total necesario para completar una response; 2. Width: registrar la cantidad y la diversity de los candidatos generados por muestreo; 3. Iteration: evaluar el beneficio marginal en accuracy de las rondas de automejora (como los trace feedback / self-correction loops).
\end{knowledgebox}

\subsection{Balance entre costo y latencia}
Chen insiste repetidamente en dar seguimiento al consumo de inference-time con métricas observables: una curva de cómputo de 1,000 dólares resulta impresionante, pero muchos escenarios de despliegue reales solo pueden asumir un presupuesto de tokens de 20 a 50 dólares. La figura \ref{fig:cost-latency}, tomada de la explicación del minuto 00:10, señala que existe un tradeoff evidente entre costo y latency, y sugiere que en los canales de baja latencia conviene priorizar el ahorro en search y verifier, mientras que en las tareas de alto valor se puede ampliar el presupuesto de tokens.

\begin{figure}[H]
\centering
\includegraphics[width=0.92\textwidth]{frame-cost-latency.jpg}
\caption{Opciones de balance entre costo y latencia en tiempo de inferencia.\protect\footnotemark}
\label{fig:cost-latency}
\end{figure}
\footnotetext{Intervalo de tiempo en el video: 00:09:40--00:10:30, muestra la escalera de presupuestos claramente distintos de $20 a $1,000.}

\begin{warningbox}{Expansión incontrolada del presupuesto}
Sin un verifier o un control de búsqueda, ampliar el presupuesto de tokens por sí solo tiende a producir una «expansión del presupuesto»: el costo se dispara sin que la accuracy mejore de forma correspondiente. Se recomienda usar caché, un conjunto de muestreo y una etapa de verifier para limitar los candidatos más costosos a las tareas de valor alto ya identificadas.
\end{warningbox}

\subsection{Programación por capas del cómputo en tiempo de inferencia}
El cómputo en tiempo de inferencia debe dividirse en ventanas de presupuesto claramente delimitadas entre las tres etapas de prompt/sampling, search e iteration. El procedimiento que Chen presenta en clase consiste en usar primero la señal de task difficulty para fijar un token cap total (por ejemplo, entre 400 y 600 tokens), y después derivar en sentido inverso la participación de cada etapa a partir de la search depth y la meta de iteration: reservar entre el 25\% y el 100 token para el prompt, ampliar la etapa de search según el número de branch, dedicar entre 80 y 120 tokens a cada ronda de iteration, y reservar además un 20\% como margen para el scoring del verifier. Este mecanismo por capas permite que, con el mismo budget, la reallocation del presupuesto entre las distintas etapas produzca estrategias diferentes que privilegien la accuracy o la latency.

\begin{table}[H]
\centering
\caption{Ejemplo de programación por capas del cómputo en tiempo de inferencia}
\begin{tabular}{p{0.25\textwidth}p{0.22\textwidth}p{0.20\textwidth}p{0.28\textwidth}}
\toprule
Etapa & Indicador clave & Presupuesto típico & Sugerencia de ingeniería \\
\midrule
Prompt / Sampling & Input tokens + example CoT & 25\%~35\% of total & Incrustar (embed) la difficulty metadata y las instrumentation hints en el prompt, para garantizar que el sampling pueda arrancar con high temperature \\
Search (branching) & Width, stepwise score, branching ratio & 40\%~50\% & Usar verifier gating para controlar la depth, de modo que cada expansión aporte un score margin \\
Iteration / Feedback & Trace iterations, consistency votes & 20\%~30\% & En la etapa de iteración es común el self-correction o el verifier replay, que consume tokens pero corrige errores de forma notable \\
Verifier / Logging & Score margin, drift detect & 10\%~15\% & Instrumentación en paralelo, como sampler metrics y verifier distribution, para mantener la latency ajustada \\
\bottomrule
\end{tabular}
\end{table}

\begin{knowledgebox}{Bucle de control de la programación por capas}
1. Fijar primero el token cap y el latency target; 2. Asignar de forma aproximada la proporción por stage y registrarla en el dashboard; 3. Observar la proporción de cada stage por query durante la ejecución y, si el drift de search/iteration es demasiado alto, corregirlo con un dial-back del prompt.
\end{knowledgebox}

\subsection{Latency-Controlled Prompting}
A nivel de Operations, para evitar que la latency se salga de control, Chen recomienda ajustar el prompt como «budget-aware templates»: indicar en el prompt «si el presupuesto supera los 400 tokens, retroceder» o «completar el razonamiento en 1 a 2 iteration». Al mismo tiempo, podemos insertar en el prompt una instrumentation hint (como `Sampling policy: alt-branch`) para que el sistema de dispatch limite el número de candidatos desde la etapa de sampling. Estas micro-guidelines ayudan a exponer el tradeoff de cost/latency directamente en el template de inference-time, en lugar de dejarlo a alertas de monitoreo after-the-fact.

\begin{warningbox}{Puntos clave del diseño del latency guard}
Los guard rail en el prompt deben ser concretos: escribir «máximo 2 iteration» es más eficaz que decir simplemente «controlar la latency»; también se puede usar `If token > 400 then reply "budget exceeded"` como límite estricto, para obligar al modelo a reportar de antemano el token usage.
\end{warningbox}

\subsection{Resumen de la sección}
La programación del cómputo en tiempo de inferencia exige medir a la vez el presupuesto de tokens, el candidate width y la iteration depth, e incorporar desde el diseño inicial del budgeting tanto el latency guard como la programación por capas del cómputo en varias etapas. Solo al tratar estos tres elementos como los ejes de la ingeniería es posible lograr una inferencia confiable y eficiente en los distintos escenarios de aplicación.
\section{Búsqueda y selección de múltiples candidatos}

\subsection{Muestreo autoconsistente y candidatos múltiples}
En el nivel de search, la autoconsistencia (Self-Consistency) se convierte en el arma central: permite que el modelo genere varios candidatos y después usa el voto por mayoría o el consistency score para elegir la final answer más estable. Como se aprecia en la figura \ref{fig:self-consistency}, a medida que el número de candidatos sube de 10 a 40, la accuracy sigue mejorando, muy por encima del ordenamiento basado en log-prob.

\begin{figure}[H]
\centering
\includegraphics[width=0.92\textwidth]{frame-self-consistency.jpg}
\caption{Curvas de Self-Consistency que comparan el ordenamiento log-prob con otras estrategias de sampling.\protect\footnotemark}
\label{fig:self-consistency}
\end{figure}
\footnotetext{Intervalo de tiempo en el video: 00:45:05--00:50:05; en la figura, el muestreo diversificado y la votación superan claramente al ordenamiento log-prob único.}

\begin{knowledgebox}{Tres elementos del muestreo con múltiples candidatos}
1. Diversidad: una temperature alta, el nucleus sampling y la perturbación del orden del prompt garantizan la diferencia entre los candidate; 2. Estrategia de selección: el voto por mayoría es mejor que el log-prob, salvo que se haya entrenado un verifier fuerte para cierta tarea; 3. Asignación de recursos: el multi-sample solo se activa en escenarios que necesitan una exactitud más alta y toleran la latencia.
\end{knowledgebox}

\subsection{Sinergia entre Self-Consistency y Search}
Self-Consistency le da a search un width más amplio, mientras que beam search/Tree-of-Thought controla el depth mediante stepwise scoring. Chen señala en el minuto 47:02 que el flujo ideal es primero usar sampling para generar entre 20 y 40 candidate, después usar search + verifier para decidir qué ruta puede seguir profundizándose y, al final, usar la votación de consistency para estabilizar la answer final. Usar simplemente el majority vote es confiable, pero para controlar el costo en inference-time todavía es necesario coordinarse con search: por un lado se garantiza la diversidad de los candidatos y, por otro, se usa el verifier para cortar los bad paths desde early stage.

\begin{warningbox}{No depender solo de uno de los módulos}
Si solo se ejecuta Self-Consistency sin agregar ningún search, se puede agotar el costo de token sin dejar de repetir errores similares; a la inversa, si solo se ejecuta search sin agregar sampling, la mayoría de las veces solo se alcanza un óptimo local. Por eso es necesario encadenar, dentro del inference-time pipeline, la generación de sample, la expansión de search y la votación de consistency.
\end{warningbox}

\subsection{Verificadores y líneas base de ordenamiento}
Si se entrena un stepwise verifier, se pueden filtrar las ramas de baja calidad desde el inicio de la generación; el trace de AlphaCode de DeepMind es un ejemplo de esto, pues documenta por sí misma la code execution y otorga una puntuación en cada paso, lo cual influye en el search posterior. Un verifier escalar debe equilibrar la estabilidad y la interpretabilidad, de lo contrario puede juzgar erróneamente una respuesta correcta y compleja como una «caja negra no interpretable».

\begin{warningbox}{La trampa de un verifier débil}
Si el verifier se inclina en exceso hacia las salidas cortas o hacia las plantillas del conjunto de entrenamiento, calificará el razonamiento de varios pasos como «puntuación baja» y lo descartará, lo cual hace que self-consistency converja hacia un conjunto de soluciones comunes pero incorrectas. El verifier debe ser sensible tanto a la accuracy como a la reasoning depth.
\end{warningbox}

\subsection{Pistas de ingeniería para la búsqueda}
En la diapositiva del minuto 52:50, Chen cuantifica aún más la salida de sampling, search y verifier: métricas como el conteo de muestreo, candidate score y voter margin deben revisarse continuamente en los inference-time logs. Un pipeline común es generar 30 response con high-temperature sampling, asignar una puntuación a cada response con un stepwise verifier y después usar majority vote para coronar la solución final. Cada etapa debe tener su propia tabla de monitoreo, de modo que en un despliegue real se pueda detectar el warning de «falta de diversity en el muestreo de hoy» o «verifier drift».

\begin{knowledgebox}{Métricas del pipeline de búsqueda}
1. Sampling depth: cantidad de candidatos y diversity; 2. Verifier quality: tendencia y estabilidad del stepwise score; 3. Voting stability: voter margin o consistency rate, usados para decidir si se necesita search adicional.
\end{knowledgebox}

\subsection{Resumen de la sección}
La autoconsistencia junto con un muestreo diversificado permite que el modelo se recupere de mistakes, pero debe combinarse con un ordenamiento por verifier de alta calidad. Todo el mecanismo requiere equilibrar accuracy, token budget y latency.

\section{Profundización en el diseño del verificador}

\subsection{Composición del stepwise verifier}
Un verifier de alta calidad se compone de al menos tres submódulos: step scoring, trace evaluation y metadata gating. Chen enfatiza que no hay que tratar al verifier como un evaluador de caja negra post-hoc, sino que hay que escribirlo dentro del pipeline de sampling/search. Por ejemplo, en una rama autogenerada de Trace, se puede hacer que el verifier asigne primero un depth score (la longitud y complejidad del chain-of-thought), luego un validity score (si cumple con el problem constraint) y, por último, un confidence margin (el margin del voto mayoritario). Estos tres indicadores constituyen la tupla de salida del stepwise verifier, que se usa como pruning gate dentro del search loop.

\begin{table}[H]
\centering
\caption{Salida del stepwise verifier}
\begin{tabular}{p{0.28\textwidth}p{0.25\textwidth}p{0.35\textwidth}}
\toprule
Indicador & Forma de cálculo típica & Función \\
\midrule
Depth score & Prompt override / reasoning depth heuristic & Evita que el verifier favorezca en exceso los outputs cortos \\
Validity score & Constraint check + domain rules & Garantiza que la branch no se desvíe de la template \\
Confidence margin & Consistency across candidates & Provee la base del margin para el majority vote \\
\bottomrule
\end{tabular}
\end{table}

\begin{knowledgebox}{Estrategia de entrenamiento del verifier}
1. Entrenar el validity + depth scorer con datos de supervised trace; 2. en inference-time, asignarle una higher probability al control token `{@validator}`; 3. combinar hashing checkpoints (como el trace hash) para evitar que el verifier haga drift al repetir steps.
\end{knowledgebox}

\subsection{Verifier gating loops}
El verifier necesita formar un ciclo cerrado con el search: cada vez que se abre y evalúa un candidate, el verifier entrega un margin y lo usa como branching priority. El gating loop que Chen muestra en el minuto 36:10 es: candidate → verifier → score queue → branch expansion/termination. Este proceso puede verse como un scheduler de múltiples colas, en el que las branches con mayor margin se empujan en paralelo hacia el tree, mientras que las de menor margin reciben early stop. Si en alguna etapa el margin cae de manera global, eso indica que la diversity del sampling podría estar en collapse, y hay que retroceder de inmediato a la template «prompt plus trace» y activar `trace replay`.

\begin{warningbox}{Riesgo de margin collapse}
Cuando el margin del verifier cae de manera global, no hay que agregar tokens de inmediato, sino revisar primero la sampling diversity y la prompt difficulty. Un margin bajo puede deberse a un dataset shift o a un prompt drift; conviene ejecutar primero `trace replay` para ver si se trata de self-correction drift o de prompt drift, y solo después decidir si se amplía el search.
\end{warningbox}

\subsection{Multi-verifier orchestration}
En la práctica se despliegan varios verifiers al mismo tiempo: un heuristic scorer rápido (por ejemplo, tf-idf + regex), un trace evaluator basado en LM y un deterministic checker (como un proof verifier). Dentro del inference-time pipeline se encadenan estas tres capas: primero el fast verifier hace un trim rápido y después entrega la high-scoring branch al slow-but-robust verifier para que la reevalúe. Chen enfatiza que el audit log debe registrar las decisions de cada capa de verifier, de modo que, ante una anomaly, se pueda ubicar «qué capa descartó al candidate».

\begin{knowledgebox}{Ritmo de coordinación de los múltiples verifiers}
1. Fast verifier: bajo cost, se usa para el early pruning; 2. LM verifier: latency media, se encarga de la reasoning depth; 3. Deterministic checker: alto costo pero alta confiabilidad, se usa con frecuencia al final de una high-stakes query; 4. registrar el output de cada capa en el dashboard, para que el equipo de observability pueda identificar qué capa provocó el margin drift.
\end{knowledgebox}

\subsection{Resumen de la sección}
El diseño del verifier no solo requiere un mecanismo de scoring de alta calidad, sino que también debe formar un budgeting gating junto con el search, e incluso proteger las high-value queries mediante multi-verifier orchestration. Solo al convertir al verifier en un inference-time control loop se puede evitar que un budget race o un drift provoquen el colapso del reasoning.

\section{Exploración en árbol y evaluación a nivel de árbol}

\subsection{Tree-of-Thought y búsqueda por ramas}
Tree-of-Thought permite abrir varias ramas en cada paso, y usa stepwise scoring para decidir qué branch vale la pena expandir. Como se muestra en la figura \ref{fig:tree-search}, beam search solo conserva los top-k, mientras que Tree-of-Thought puede filtrar la promisingness ya en el step-level.

\begin{figure}[H]
\centering
\includegraphics[width=0.92\textwidth]{frame-tree-search.jpg}
\caption{Esquema comparativo entre Tree-of-Thought, beam search y autoconsistencia. \protect\footnotemark}
\label{fig:tree-search}
\end{figure}
\footnotetext{Intervalo de tiempo en el video: 01:01:24--01:04:30, muestra la posición de step evaluation dentro de la búsqueda en árbol.}

\begin{importantbox}{Tres perspectivas prácticas de la búsqueda en árbol}
1. No todas las ramas deben expandirse; el orden de prioridad puede asignarlo directamente el step verifier; 2. step-level scoring puede descartar rutas de baja calidad antes de que el branch termine, lo que reduce el costo de tokens; 3. beam search puede combinarse con sampling: primero se usa sampling para generar diverse seed, y luego beam/tracing para filtrar promising ones.
\end{importantbox}

\begin{knowledgebox}{Lógica operativa del control en árbol}
En Tree-of-Thought, el control loop suele usar tres pasos: «branch proposal → step scoring → prune»: el primer paso es hacer warm-start de varios branch, el segundo usa un verifier para dar una puntuación a cada branch, y el tercero poda según el budget. Chen recomienda registrar en el training log el token count y el score de cada branch, para poder identificar «en qué step el score drift es más marcado».
\end{knowledgebox}

\begin{warningbox}{El branching en árbol no puede expandirse sin límite}
Aunque tree search puede mejorar el accuracy en early stage, si no hay una estrategia de pruning clara, el número de branch crecerá de forma exponencial. Es indispensable establecer un token threshold (por ejemplo, explorar como máximo 3 niveles), incorporar stepwise verifier gating y usar logs para monitorear el branching ratio.
\end{warningbox}

\subsection{Lazo de control en árbol y autoadaptación del presupuesto}
En el pipeline de Tree-of-Thought debe existir un lazo de control que decida cuándo expandir, cuándo hacer prune y cuándo hacer roll back. La lógica de control que Chen presentó en vivo es: al terminar cada ronda de step se revisa si `branch score / step token` supera el threshold; si está por debajo del threshold, el branch se envía de inmediato a la aborted queue y los token restantes se recuperan para el global budget. Este loop, mediante SLO instrumentation (como `branching\_ratio` y `score\_per\_token`), expone las ramas inválidas de tree search que pasaban inadvertidas.

\begin{table}[H]
\centering
\caption{Indicadores de control de tree search}
\begin{tabular}{p{0.32\textwidth}p{0.25\textwidth}p{0.32\textwidth}}
\toprule
Indicador & Significado del monitoreo & Pipeline action \\
\midrule
Branch score & Determina la promisingness de la rama & Cuando Score < 0.2, se hace prune y se dispara `trace replay` \\
Branching ratio & Observa la relación entre active branches y token & Cuando > 28\% se limitan los new branches a 2 \\
Score per token & Controla depth y cost & Una caída indica search drift; hay que bajar la temperature o hacer early stop \\
\bottomrule
\end{tabular}
\end{table}

\begin{importantbox}{Lecciones del control loop de Tree}
1. Se establece `early score check`, es decir, calcular el partial score antes de que el branch termine; 2. se usa per-branch token cap para obligar a que las ramas no se extiendan sin límite; 3. cuando el control loop presenta una anomalía (por ejemplo, cuando branching ratio se dispara), se activa de inmediato conservative mode (bajar el prompt temperature + limitar iteration).
\end{importantbox}

\subsection{Resumen de la sección}
Estructuras en árbol como Tree-of-Thought llenan, dentro de inference-time search, el gap entre beam search y sampling; basta con combinarlas con un step verifier sólido y un lazo de control para lograr, con un budget limitado, high-quality breadth exploration.
\section{Retroalimentación y automejora iterativa}

\subsection{Trace feedback y evaluación por pasos}
Chen presenta el trace feedback de AlphaCode, que usa lenguaje natural para reify la traza de ejecución del modelo, formando un stepwise verifier. Al hacer que el modelo genere su propio line-by-line debug, este puede identificar qué estados intermedios son promising, y usa una evaluación text-level que influye en el search posterior.

\begin{figure}[H]
\centering
\includegraphics[width=0.92\textwidth]{frame-trace-feedback.jpg}
\caption{Captura de ejemplo de trace feedback, que muestra la visualización de la line-by-line explanation.\protect\footnotemark}
\label{fig:trace-feedback}
\end{figure}
\footnotetext{Intervalo de tiempo en el video: 01:11:20--01:12:15, donde se explica el funcionamiento del trace feedback.}

\begin{importantbox}{El valor de más alto nivel del trace feedback}
1. Convierte la step evaluation en un natural language trace, haciendo que el propio modelo se convierta en verifier; 2. AlphaCode observa la context renovation a través del trace, lo que le permite detectar subtle bugs durante la etapa de code generation; 3. Un trace más rico puede usarse directamente como una signal de recompensa de RL, lo que ofrece una vía de mejora en inference-time para sistemas como AlphaProof y AlphaGeometry.
\end{importantbox}

\subsection{Prompt de autocorrección y riesgos}
Cuando no hay un oracle verifier, la autocorrección a menudo transforma un error en otro error distinto; Chen advierte de esto usando un ejemplo negativo en 01:13:23. Al diseñar un general-purpose feedback prompt, es necesario resaltar la consistency, incorporar el trace context y establecer guard rails del tipo «no alterar las partes ya confirmadas».

\begin{knowledgebox}{Direcciones para ajustar el prompt de autocorrección}
1. Resaltar la consistency: indicarle al modelo que «enumere dos posibles pasos nuevos», en lugar de una guía ambigua como «revisa la respuesta»; 2. Inyectar trace context: usar de nuevo la line-by-line explanation como prompt, para ayudar al modelo a juzgar la correctness con su propio reasoning; 3. Controlar los saltos: agregar la instrucción «conservar el razonamiento ya confirmado» para evitar que el modelo, por feedback drift, revierta las partes ya verificadas.
\end{knowledgebox}

\begin{warningbox}{Desventajas del self-correction sin oracle}
Cuando solo hay general feedback y no hay ground truth, el modelo, al no poder distinguir con claridad qué response es confiable, termina reescribiendo de forma conservadora una respuesta correcta o abandonando el reasoning de varios pasos, lo que provoca que la accuracy disminuya en vez de mejorar.
\end{warningbox}

\subsection{Resumen de la sección}
El trace feedback y el self-correction necesitan un reliable verifier y guard rails en el prompt; de lo contrario, el modelo retrocede después de varios rounds. Al diseñar el feedback hay que atender a la vez a la guía de consistency, al trace context y a los límites de interpretabilidad.
\section{Fusión de estrategias y programación}

\subsection{Combinación de prompt, búsqueda y trace}
Chen enfatiza que las estrategias anteriores no están aisladas: se puede usar primero Tree-of-Thought para expandir las branch, luego usar votación por self-consistency para estabilizar la final answer, y por último dejar que el trace feedback itere para verificar. En otras palabras, el pipeline tiene la forma prompt → search → iteration, y cada eslabón tiene su propia métrica y su propio budget.

\subsection{Programación de presupuesto y métricas}
El despliegue real necesita programar un triple presupuesto de token, branch e iteration. En una tarea compleja se puede asignar más token en las etapas iniciales, usar un verifier para controlar el número de branch en la etapa de search, y limitar los feedback rounds en la etapa de iteration. En cuanto a métricas, hay que registrar a la vez accuracy, token count y latency.

\begin{knowledgebox}{Principios de programación de presupuesto}
1. Escalamiento automático: usar la task difficulty signal para ajustar el token/branch depth; 2. cola de prioridad: apoyarse en el verifier para puntuar cada branch y controlar el orden de expansión; 3. representación observable: registrar el F1, el token count y la latency de cada módulo (prompt, muestreo, trace) durante el deployment, para poder reponderar los recursos en cualquier momento.
\end{knowledgebox}

\begin{importantbox}{Consejo práctico para la fusión de estrategias}
No traten a prompt, search y trace como sistemas independientes, sino como una cadena de presupuesto: cuando el search depth aumenta, hay que descontar automáticamente el token budget o limitar el iteration; cuando la profundidad de iteration aumenta, hay que regresar el search branching a un valor seguro. Solo así se puede evitar, en el despliegue, el fenómeno de que «todos los módulos se extienden por separado y el presupuesto se expande exponencialmente».
\end{importantbox}

\subsection{Resumen de la sección}
La fusión de estrategias significa tratar el prompt, la búsqueda y el trace como un pipeline cíclico cerrado. Mediante la programación de presupuesto, la cola de prioridad y el ciclo de métricas, se puede mantener una alta precisión y una alta interpretabilidad bajo inference-time compute.
\section{Interpretación de Benchmark y direcciones de investigación}

\subsection{Interpretación de las curvas de Arc AGI / O3}
Arc AGI, como indicador central de un reasoning benchmark, permitió inicialmente que un presupuesto de 1,000 dólares alcanzara 87.5\% de accuracy, mientras que por debajo de 20 dólares solo se mantenía en 25\%. Chen señala que esta curva no es simplemente cuestión de «pagar para lograrlo», sino que requiere coordinar al mismo tiempo sampler, search y verifier, para que cada partida de presupuesto se convierta en accuracy efectiva. La práctica de demos como OpenAI O3, AlphaProof/AlphaGeometry, depende firmemente del inference-time compute, y no del dominio de entrenamiento.

\begin{table}[H]
\centering
\caption{Comparación entre presupuesto y desempeño en tiempo de inferencia de los principales Benchmark}
\begin{tabular}{p{0.30\textwidth}p{0.28\textwidth}p{0.32\textwidth}}
\toprule
Benchmark & Token / Cost budget & Observed behavior \\
\midrule
Arc AGI & \$20~\$1,000 per query & 25\% → 87.5\% accuracy; the curve saturates after \$500 without better search/verifier \\
OpenAI O3 & \$20~\$200 & Consistency + trace feedback gives human-level accuracy once budget crosses \$100 \\
AlphaProof / AlphaGeometry & 50~200 token steps & RL-based reward tweaks inference-time policy to maintain correctness under limited latency \\
\bottomrule
\end{tabular}
\end{table}

\subsection{Tablero de Benchmark y niveles de presupuesto}
Chen recomienda construir para cada benchmark un tablero de ``budget bin'': dividir el presupuesto de token/cost en tres niveles, Conservative, Balanced y Aggressive, que corresponden respectivamente a los casos de `20~80 token`, `80~300 token` y `>300 token`. Cada bin debe registrar accuracy, stepwise verifier margin y token per vote, para poder ubicar rápidamente en tiempo de ejecución «en qué bin se encuentra el estado actual». Esta estratificación permite que el equipo, al publicar los resultados del benchmark hacia afuera, diga directamente «en el bin Balanced, nuestra accuracy es de 87.5\% pero la latency es de 3.4s».

\begin{table}[H]
\centering
\caption{Ejemplo de tablero de Benchmark}
\begin{tabular}{p{0.26\textwidth}p{0.30\textwidth}p{0.30\textwidth}}
\toprule
Bin & Configuración típica & Indicadores de monitoreo \\
\midrule
Conservative & $\leq 80$ token, $\leq 2$ iterations & Accuracy, sampling diversity, latency \\
Balanced & 80~300 token, 2~3 iterations & Verifier margin, branching ratio, token per candidate \\
Aggressive & > 300 token, 3+ iterations & Score per token, trace iterations, SLO breach rate \\
\bottomrule
\end{tabular}
\end{table}

\begin{importantbox}{Recomendaciones para implementar el tablero de Benchmark}
1. Configurar para cada bin un `alert level`, y cambiar de bin cuando la latency supere el target o el verifier margin caiga por debajo del threshold; 2. hacer push del state del bin al dashboard, para que el equipo de operaciones pueda leer rápidamente «en qué nivel está corriendo el benchmark hoy»; 3. combinar esto con el monitoreo de score-per-token dentro del bin, para detectar a tiempo el early warning de que el token budget se está saliendo de control.
\end{importantbox}

\subsection{Preguntas de investigación abiertas}
\begin{itemize}
\item ¿Cómo ampliar el search depth y el iteration depth con compute limitado en escenarios de baja latency?
\item ¿Es posible diseñar un «budget-aware verifier» que reporte en tiempo real el tradeoff entre token y candidate, sin esperar a que termine la inference?
\item ¿Cómo garantizar la generalidad del trace feedback en el reasoning multimodal o de imagen más texto?
\item ¿Cómo, en un pipeline de despliegue real, encadenar indicadores como sampling diversity, stepwise scaling metric y trace drift en un SLO ejecutable?
\end{itemize}

\begin{itemize}
\item ¿Puede un pipeline de inferencia ajustar de manera adaptativa el branching ratio y el iteration depth mediante learning-to-rank, en lugar de fijar umbrales manualmente?
\item Para un adversarial prompt, ¿cómo combinar el verifier margin con el prompt-editing loop, para evitar que un razonamiento erróneo se refuerce mediante el feedback?
\item En el reasoning multimodal, ¿qué incertidumbres persisten en la correspondencia entre el presupuesto de token y el vision token length?
\end{itemize}

Estas preguntas se mencionaron repetidamente en el Q\&A posterior a las 01:19:00, y Chen sugiere explorarlas mediante el enfoque «metric-first hypothesis»: primero proponer una metric central (como `score per token`) y luego iterar el ajuste del pipeline, en lugar de ampliar el token budget guiándose por la intuición.

\begin{knowledgebox}{Asociaciones de direcciones de investigación}
Las preguntas abiertas pueden abordarse desde tres ángulos: budgeting instrumentation, verifier robustness y la ciencia del trace. Chen recomienda escribir los inference-time logs (sampling count/delay/score) en el dashboard, para proporcionar los datos base que permitan en el futuro ajustar automáticamente la estrategia.
\end{knowledgebox}

\subsection{Resumen de la sección}
A través de la interpretación de las curvas de budget/accuracy de Arc AGI, O3 y la serie de DeepMind, vemos que el inference-time compute es a la vez un parámetro de ingeniería y un detonante de direcciones de investigación. En el futuro será necesario mantener una high accuracy bajo low-latency, e integrar los indicadores de trace/verifier/gating en el pipeline de despliegue.

\section{Despliegue y observación}

\subsection{Monitoreo de métricas inference-time}
Chen ve el inference-time pipeline como una «caja negra»: tras recibir el prompt/block de entrada, pasa sucesivamente por sampling, search y trace, y cada etapa debe generar metrics. Las metrics típicas incluyen sampling depth, branching ratio, verifier score distribution, trace iteration count y latency. Se recomienda graficar estas curvas en el dashboard agrupadas cada 5 minutos; en un sistema real, en cuanto el sampling depth caiga de manera repentina o el trace iteration aumente bruscamente, debe activarse una investigación.

\begin{knowledgebox}{Cadena de herramientas de observabilidad}
1. Sampling log: registra prompt + generated candidates + token count; 2. Search log: registra beam/tree status (depth, score, branch id); 3. Trace log: indica en cada ronda de self-correction si hay un error pattern nombrado por humanos; 4. Alert rules: configura un collector que emita una alerta cuando la diversidad de sample sea inferior al baseline o el trace drift supere el threshold.
\end{knowledgebox}

\subsection{Budget SLO y alertas}
Al desplegar el inference-time compute pipeline, es común usar una tabla de SLO que enumere los límites superiores de token budget, search depth e iteration depth, junto con alertas del tipo «cuando average token per query > budget × 1.1» o «search branching > 30%». Chen recomienda agregar una etapa de «double-check» en el SLO, y usar los dos score del stepwise verifier para asegurar que una high-stakes query no falle debido a un sampling drift.

\begin{table}[H]
\centering
\caption{Ejemplo de SLO de inference-time}
\begin{tabular}{p{0.32\textwidth}p{0.32\textwidth}p{0.32\textwidth}}
\toprule
Metric & Target & Alert condition \\
\midrule
Token budget per query & $\leq 400$ tokens & 5-minute moving average > 440 tokens \\
Sampling diversity & 80\% unique candidates & 3 consecutive minutes diversity < 60\% \\
Iteration depth & $\leq 2$ feedback rounds & iteration count > 3 for top 1\% queries \\
Verifier score margin & $\geq 0.15$ & margin drop > 0.05 for 100 queries \\
\bottomrule
\end{tabular}
\end{table}

\begin{warningbox}{Recomendaciones para el ajuste de alertas}
No permitas que cada trace drift dispare una alerta, pues eso provocaría alarm fatigue. La experiencia de Chen es fijar el umbral de falsos positivos un poco más alto, y programar una investigación manual solo cuando exista un problema real, apoyándose en la función de «replay» para que el equipo revise el sampling log.
\end{warningbox}

\subsection{Resumen de la sección}
Al desplegar el inference-time pipeline es necesario que observabilidad, SLO, alertas y log replay trabajen de manera coordinada, para poder mantener accuracy y reliability en la capa de production, y al mismo tiempo mantener el budget dentro de un rango controlable.

\section{Ejemplos prácticos y plantillas de instrucciones}

\subsection{Práctica del Prompt Template}
Un inference-time prompt ideal incluye context, task difficulty, example CoT, instrumentation hints y final instruction. Por ejemplo, la siguiente plantilla se usa con frecuencia en reasoning de varios pasos:

\begin{lstlisting}[language=Python, caption={Inference-time prompt skeleton}]
System: You are a high-precision reasoning assistant.

Input context: {problem description}
Difficulty: {mark as "hard"/"critical"}
Examples:
  - Example 1 CoT: ...
  - Example 2 CoT: ...
Instruction: Think step-by-step. If uncertain, enumerate two next steps and sample multiple candidates.
\end{lstlisting}

\begin{importantbox}{Técnicas para poner en práctica la plantilla de prompt}
Vincula difficulty como metadata dentro del prompt, para que la sampling policy ajuste automáticamente temperature y token limit; usa ejemplos de casos que expresen el contraste «wrong → right», para recordarle al modelo que debe extender el reasoning en profundidad.
\end{importantbox}

\subsection{Self-Consistency Orchestration}
En la etapa de sampling, primero se generan continuamente entre 20 y 40 responses con high-temperature; en la etapa de search, estas responses se toman como seed y se activa beam/tree search para extender más las promising branches; en la etapa de iteration, se ejecutan trace feedback y la puntuación del verifier, y después se elige el answer final con base en majority vote/consistency.

\begin{knowledgebox}{Orchestration checklist}
1. Sampling: registra candidate count + diversity; 2. Search: track branch depth, stuck rate; 3. Iteration: log trace iteration count, verification score, self-correction delta.
\end{knowledgebox}

\subsection{Resumen de la sección}
Los ejemplos prácticos ayudan a concretar la idea abstract de inference-time: podemos apoyarnos en el prompt skeleton, las orchestrations de las tres etapas sampling/search/iteration, y el instrumentation checklist para escribir el pipeline como un template reutilizable.

\section{Riesgos y gobernanza}

\subsection{Identificación de riesgos}
Los risk comunes incluyen budget fuera de control (el token budget se rebasa de pronto), verifier drift (caída del score margin), trace drift (la self-correction deja de mejorar la accuracy) y sampling collapse (caída de la diversity). Cada vez que aparecen estos riesgos, Chen recomienda revertir de inmediato la sampling policy (temperatura, nucleus) y reproducir el proceso de decisión con «trace replay».

\begin{warningbox}{Señales de riesgo y respuesta inicial}
1. Token budget > 110\% of SLO → reduce search width; 2. Verifier margin drop > 0.05 → rerun sampling with new prompts; 3. Trace iteration spike → throttle iteration depth and inspect trace log.
\end{warningbox}

\subsection{Gobernanza y simulacros}
En el plano de la gobernanza, es necesario establecer un runbook de inference-time, que incluya budget gatekeeper, trace simulator y post-mortem template. Cuando el pipeline presenta anomalías, se puede usar el runbook para cambiar de inmediato a conservative mode (reducir el token budget, limitar el search depth) y registrar el log de la anomalía en el backlog, para usarse en futuras iteraciones de prompt.

\begin{table}[H]
\centering
\caption{Activos necesarios para los simulacros de gobernanza}
\begin{tabular}{p{0.32\textwidth}p{0.32\textwidth}p{0.32\textwidth}}
\toprule
Asset & Purpose & Trigger \\
\midrule
Budget gatekeeper dashboard & Muestra el uso de token/search/iteration y admite manual override & Budget SLO breach \\
Trace simulator & Replay candidate + verifier steps for debugging & Consistency drop or new failure mode \\
Post-mortem template & Capture root cause, fix, and regression test plan & Any production inference failure \\
\bottomrule
\end{tabular}
\end{table}

\subsection{Resumen de la sección}
Con una identificación de riesgos clara, un runbook y activos de gobernanza, el Inference-time pipeline puede mantenerse estable a lo largo de múltiples iteraciones y garantizar que los apartados de prompt/search/trace operen dentro del presupuesto y conforme a las métricas.
\section{Análisis de casos y siguientes pasos}

\subsection{Arc AGI Puzzle Walkthrough}
En esta curva de benchmark de Arc AGI, podemos usar los siguientes pasos para impulsar el inference loop: 1) en el prompt se declara de antemano que este puzzle es high-difficulty; 2) en la etapa de sampling se generan entre 30 y 40 candidate response; 3) en la etapa de search estos candidate se toman como tree seed, y con un stepwise scorer se calcula el promisingness además de introducir un probability threshold; 4) en la etapa de iteration se ejecuta trace feedback, se autorrevisa la reasoning chain de cada response, y con ayuda del verifier margin se orienta la elección por votación mayoritaria. Todo el proceso también registra every candidate's token cost y latency, para future tuning.

\begin{importantbox}{Observación del pipeline de Arc AGI}
Hay que considerar Arc AGI como un campo de práctica de todo el proceso: prompt → sampling → search → trace → voting, cada eslabón necesita instrumentation; si el token cost se duplica de pronto, hay que revisar si el sampling diversity y el verifier margin se han derrumbado, y luego usar el runbook para volver al conservative mode.
\end{importantbox}

\subsection{Despliegue en tiempo de inferencia para tareas de alto valor}
Para tareas competitivas de Code/Math como AlphaProof, el pipeline normalmente también agrega extra verification: 1) Run deterministic verifier (e.g. proof checker) as soon as the candidate completes; 2) If verification fails, reroute to alternative candidate with similar margin; 3) Log the mismatch and feed trace back to sampling controller for prompt iteration. Este tipo de flujo convierte el inference-time compute en un closed-loop control rather than a one-shot guess.

\begin{knowledgebox}{Puntos clave del despliegue de alto valor}
1. Usar el deterministic verifier como fallback, para evitar que el majority vote tome un wrong candidate como respuesta correcta; 2. usar trace replay para reproducir cada verification failure, de modo que se puedan generar nuevos examples mediante prompting; 3. registrar las metrics (token cost, verifier success, trace iteration) del close-loop pipeline en el performance dashboard.
\end{knowledgebox}

\subsection{Resumen de la sección}
El análisis de casos mostró cómo llevar a la práctica el compute pipeline en tiempo de inferencia en tareas del tipo Arc AGI y AlphaProof: primero se determina el nivel del prompt, luego se hace que sampling/search/trace constituyan los controles, y finalmente se escribe la retroalimentación de verification en el iteration loop para la mejora continua.

\section{Síntesis y ampliación}
Esta clase gira en torno al inference-time compute: primero se deja que el prompt libere tokens, después se usa sampling/search para ampliar el solution space y, finalmente, se corrige el resultado mediante trace + self-correction. Las tres etapas se pueden combinar según se necesite, y en tareas de reasoning de alta calidad mejoran de manera notable la accuracy y la reliability.

También repasamos el compute scheduling (programación triangular), el mecanismo de coordinación entre self-consistency/search, el loop de control de stepwise verifier y Tree-of-Thought, así como la interpretación de los datos de benchmark, para formar un mapa de ruta completo que va del prompt al deployment.

\begin{importantbox}{Puntos clave para confirmar la implementación}
1. Dar seguimiento continuo, en la ingeniería, a las tres categorías de métricas: sampling, search y trace; 2. usar la curva de benchmark (por ejemplo, Arc AGI) como referencia para el budget gating; 3. incorporar el deployment SLO, el alert y el budget instrumentation al CI/CD pipeline, para garantizar que el inference-time compute no exceda un rango controlable.
\end{importantbox}

\begin{table}[H]
\centering
\caption{Patrones centrales del cómputo en tiempo de inferencia}
\begin{tabular}{p{0.28\textwidth}p{0.32\textwidth}p{0.32\textwidth}}
\toprule
Estrategia & Beneficio clave & Costo típico \\
\midrule
Ampliar el presupuesto de CoT & Permite que el modelo se acerque a la human-level accuracy en problemas difíciles & Más tokens, más latency, mayor sensibilidad en el diseño del prompt \\
Muestreo de múltiples candidatos y self-consistency & Permite recuperarse de mistakes; útil en escenarios sin verifier & Cada candidate implica su propio compute, y requiere un voter pool y una representación de consistencia \\
Tree-of-Thought y step verifier & Eliminar bad paths desde etapas tempranas, y mejorar la interpretabilidad en conjunto con el verifier & El número de branching se dispara, y hace falta diseñar una quality function \\
Trace feedback + self-correction & La puntuación line-by-line puede mejorar la accuracy de forma progresiva & Sin oracle se produce feedback drift; hace falta prompt guard rails \\
\bottomrule
\end{tabular}
\end{table}

\subsection{Resumen de la sección}
Las técnicas de tiempo de inferencia necesitan integrar el prompt, el search y el trace en un pipeline, y mediante la programación del presupuesto y el bucle de métricas hacer que cada módulo aporte su máximo beneficio en cada etapa.

\subsection{Lecturas adicionales}
\begin{itemize}
\item Wang et al., en «Self-Consistency improves chain of thought reasoning», explican el efecto de scaling del sampling + majority vote.
\item Zhou et al., en «Tree-of-Thought», proponen combinar el search con la step evaluation mediante una exploración en forma de árbol.
\item Artículos sobre trace feedback como «Trace-Enhanced LLM Safety» explican cómo el natural language trace ayuda al verifier.
\item La serie AlphaProof/AlphaGeometry/AlphaCode de DeepMind muestra los resultados de combinar inference-time RL + search + trace en tareas complejas de reasoning.
\end{itemize}

\end{document}
